% BEGIN REV09_GERMEN_DODECAFASICO
\section{Germe dodécaphasique et récupération des branches unimodales
}
\label{corpus9:iv:germen}

La roue orientée \(U_{12}^{\mathrm{sign}}\), reçue du retour
TPK, détermine dans le secteur uniforme \(\eta=0\) les poids
\(w_m=1/12\). Nous fixons la fonctionnelle de phase et ses représentants
réels ordonnés

\begin{equation}
 \phi_m=\pi_{\mathrm{HMT}}\frac{3m-1}{18},
 \qquad m=1,\ldots,12.
 \label{corpus9:iv:germen-fases}
\end{equation}
Le relèvement réel des phases fait partie du lecteur :
les moments suivants sont calculés sur ces représentants.
La construction provient de la roue, de son orientation, de ses poids
et de sa fonctionnelle ; une modification de l’une quelconque de ces données
modifie la collection résultante.

\subsection{Normalisation et profil analytique}
Soient \(a_m=(3m-1)/18\) et
\(D_0=\frac1{12}\sum_{m=1}^{12}a_m^2\). Les sommes entières
\[
 \sum_{m=1}^{12}(3m-1)^2=5394=6\cdot899,\qquad
 \sum_{m=1}^{12}(3m-1)^4=4294614
\]
donnent \(D_0=899/648\). La condition de normalisation quadratique
\(s_0^2\pi_{\mathrm{HMT}}^2D_0=2\) fixe

\[
 s_0=\frac{36}{\pi_{\mathrm{HMT}}\sqrt{899}},
 \qquad k_m=s_0\phi_m=\frac{2(3m-1)}{\sqrt{899}}.
\]
L’échelle angulaire s’annule dans cette normalisation avant
d’itérer le profil
\begin{equation}
 u_0(x)=1-\frac1{12}\sum_{m=1}^{12}\cos(k_mx),
 \qquad f_\lambda(x)=1-\lambda u_0(x).
 \label{corpus9:iv:germen-perfil}
\end{equation}

\begin{proposition}[Représentation exacte et criticité
]
\label{corpus9:iv:germen-analitico}
Le profil \(u_0\) est une fonction entière paire et satisfait

\[
 u_0(x)=x^2-\frac{715769}{2424603}x^4+O(x^6).
\]
Avec \(y=x/\sqrt{899}\), son expression fermée est
\begin{equation}
 u_0(x)
  =1-\frac{\cos(37y)\sin(36y)}{12\sin(3y)}
  =1-\frac{\sin(73y)-\sin y}{24\sin(3y)},
 \label{corpus9:iv:germen-cerrado}
\end{equation}
où les quotients sont prolongés en leurs singularités éliminables.
Pour \(0<\lambda\leq2/u_0(1)\), \(f_\lambda\) est une
application réelle analytique et paire de \([-1,1]\) dans lui-même, avec un
unique point critique quadratique en zéro et une dérivée schwarzienne négative
hors de ce point.
\end{proposition}
\begin{proof}
La somme finie de cosinus est entière et paire. Son développement de
Taylor utilise
\[
 \frac1{12}\sum_m k_m^2=2,\qquad
 \frac1{24\cdot12}\sum_m k_m^4=\frac{715769}{2424603},
\]
obtenus à partir des deux sommes entières précédentes. Pour la forme
fermée, la progression des arguments est \((6m-2)y\). La somme
géométrique des exponentielles, suivie de la partie réelle, donne

\[
 \sum_{m=1}^{12}\cos((6m-2)y)
   =\frac{\sin(36y)}{\sin(3y)}\cos(37y).
\]
L’identité \(2\cos(37y)\sin(36y)=\sin(73y)-\sin y\)
prouve \eqref{corpus9:iv:germen-cerrado}. La somme finie définit
son prolongement en tous les zéros du dénominateur.

Pour \(0<x\leq1\), tous les arguments satisfont
\(0<k_mx\leq70/\sqrt{899}<\pi_{\mathrm{HMT}}\). Par conséquent,

\[
 u_0'(x)=\frac1{12}\sum_m k_m\sin(k_mx)>0,\qquad
 u_0'''(x)=-\frac1{12}\sum_m k_m^3\sin(k_mx)<0.
\]
La parité fournit les signes opposés dans l’intervalle
négatif et prouve l’unicité du point critique. En outre,
\(f_\lambda''(0)=-2\lambda\ne0\). Comme
\(0\leq u_0(x)\leq u_0(1)\) dans \([-1,1]\), l’intervalle
indiqué pour \(\lambda\) donne
\(-1\leq f_\lambda(x)\leq1\). Enfin,
\[
 S(f_\lambda)
  =\frac{u_0'''}{u_0'}
   -\frac32\left(\frac{u_0''}{u_0'}\right)^2<0
 \qquad(0<|x|\leq1).
\]
Le premier quotient est négatif de part et d’autre de l’origine,
et le terme restant est non positif.
\end{proof}

\subsection{Inversion par branche et mémoire de l’itinéraire
}
\begin{proposition}[Récupération du pliage
]
\label{corpus9:iv:germen-inversa}
Soient \(\lambda>0\) et
\(v=(u_0|_{[0,1]})^{-1}\). Pour
\(y\in[1-\lambda u_0(1),1]\), l’antécédent dans une branche est
\begin{equation}
 x=\sigma\,v\!\left(\frac{1-y}{\lambda}\right),
 \qquad \sigma=\operatorname{sgn}x.
 \label{corpus9:iv:germen-rama}
\end{equation}
En \(y=1\), la préimage est nulle et on fixe \(\sigma=0\).

\end{proposition}
\begin{proof}
La continuité et la monotonie stricte de \(u_0\) sur
\([0,1]\) fournissent l’inverse \(v\).
L’équation \(y=1-\lambda u_0(x)\), avec la parité,
détermine \(|x|=v((1-y)/\lambda)\) ; le signe détermine
l’orientation. Pour \(y=1\), l’unicité du minimum de \(u_0\)
donne \(x=0\). Hors du point critique, l’inverse est différentiable.
Sa dérivée devient singulière en arrivant à l’origine, où
\(u_0(x)=x^2+O(x^4)\).

\end{proof}

Pour une distribution finie sur des paires \(\pm x_i\), avec
\(x_i>0\) distincts, soient \(p_i^\pm\) leurs probabilités,
\(p_i=p_i^++p_i^-\), \(Y=f_\lambda(X)\) et
\(\Sigma=\operatorname{sgn}X\). On admet également une masse
en zéro, dont le terme d’incertitude de branche est nul. Avec
des logarithmes naturels,
\begin{align}
 H(X\mid Y,\Sigma)&=0,\\
 H(X\mid Y)&=H(\Sigma\mid Y)
   =\sum_{i:p_i>0}p_i\,h(p_i^+/p_i),\\
 h(p)&=-p\log p-(1-p)\log(1-p).
 \label{corpus9:iv:germen-entropia}
\end{align}
La preuve consiste à regrouper la distribution selon les fibres
de \(f_\lambda\). La monotonie distingue les modules \(x_i\),
et \eqref{corpus9:iv:germen-rama} distingue les deux signes
à l’intérieur de chaque fibre. Avec des probabilités égales dans chaque paire
et sans masse en zéro, l’incertitude vaut \(\log2\).

Si \(0<\lambda\leq2/u_0(1)\), la seconde itération est
définie sur le même intervalle et vérifie
\[
 H(X\mid f_\lambda^2(X))
 =H(X\mid f_\lambda(X))
  +H(f_\lambda(X)\mid f_\lambda^2(X)).
\]
En effet, les deux applications intermédiaires sont déterministes,
et la règle de la chaîne pour l’entropie de
\((X,f_\lambda(X))\), conditionnellement à \(f_\lambda^2(X)\),
donne l’identité. La récupération successive utilise les signes
de chaque pas. Pour \(\lambda=0\), l’application est constante
et exige de conserver en outre le module de l’entrée.

\subsection{Domaine du retour remis à l’échelle}
Pour \(a=-f(1)\ne0\), soit \(h_a(x)=-ax\). L’opérateur
\[
 \mathcal R(f)=h_a^{-1}\circ f^2\circ h_a
\]
est un retour remis à l’échelle sur \([-1,1]\) lorsque l’intervalle
\(J=h_a([-1,1])\) appartient au domaine de \(f^2\) et satisfait
\(f^2(J)\subseteq J\). L’échelle \(h_a\) est bijective ;
le retour en deux pas conserve son interprétation avec
l’itinéraire et l’application \(f\). La reconstruction d’une
racine itérative à partir de \(\mathcal R(f)\) et \(a\) constitue
une condition supplémentaire. La proposition
\ref{corpus9:iv:germen-analitico} établit le germe et sa
classe critique. L’existence et l’unicité d’un point fixe du
renormalisateur, sa décomposition spectrale stable et instable
et la transversalité de cette collection sont les hypothèses
supplémentaires du passage à une limite universelle de renormalisation.
% END REV09_GERMEN_DODECAFASICO
